\documentclass[12pt]{amsart}
\usepackage[T1]{fontenc}
\usepackage[utf8]{inputenc}
\usepackage{amssymb}
\usepackage{float}
\usepackage[backend=biber, style=numeric,  bibencoding=utf8]{biblatex}
\addbibresource{infinite_tetrads.bib}

\newtheorem{thm}{Theorem}[section]
\newtheorem{lem}[thm]{Lemma}

\theoremstyle{remark}

\newtheorem{remark}[thm]{Remark}



\begin{document}

\title{Infinite tetrads of congruent factorials}


\author{Octavio A.  Agustín-Aquino}
\address{Instituto de Física y Matemáticas, Universidad Tecnológica de la Mixteca, Huajuapan de León, Oaxaca, México}
\email{octavioalberto@mixteco.utm.mx}
\urladdr{www.utm.mx/~octavioalberto} 

\author{José Hernández Santiago}
\address{Escuela Superior de Matemáticas No. 2, Universidad Autónoma de Guerrero, Ciudad Altamirano, Guerrero, México}
\email{15620@uagro.mx}
\urladdr{https://euclid.uagro.mx/}
\subjclass[2020]{11B65, 11A07, 11A41}

\begin{abstract}
We prove that there exist infinitely many primes $p$ such that there exist $k_{1},k_{2},k_{3},k_{4}$ distinct
such that
\[
k_{1}! \equiv k_{2}! \equiv k_{3}! \equiv k_{4}! \equiv 1 \pmod{p}.
\]
\end{abstract}


 \maketitle


\section{Introduction}

Following \cite{GKP94}, we write $a\backslash b$ to mean $a$ divides $b$. The following fact is relatively well-known.

\begin{lem}[Wilson's symmetry]
Let $p$ be an odd prime and $0\leq n \leq p-1$ an odd integer such that $n!\equiv 1\pmod{p}$. Then
\[
(p-1-n)!\equiv 1 \pmod{p}.
\]
%and, in particular,
%\[
%(p-2)! \equiv 1 \pmod{p}.
%\]
\end{lem}
\begin{proof}
By Wilson's theorem we have $(p-1)!\equiv -1 \pmod{p}$. Then
\[
 (p-1)! = (p-1-n)! \prod_{k=1}^{n}(p-k).
\]

Thus, modulo $p$, we have
\begin{align*}
 (p-1-n)! \prod_{k=1}^{n}(p-k) &\equiv (p-1-n)!(-1)^{n}n!\\
 &\equiv -(p-1-n)! \equiv -1 \pmod{p},
\end{align*}
hence $(p-1-n)!\equiv 1 \pmod{p}$. 
\end{proof}

Wilson's symmetry is crucial for the following construction.

\begin{lem}[Valid tetrad]
Let $q\geq 3$ be an odd integer, and $p$ a prime divisor of $q!-1$. If $p$ is distinct from
$q+2$ and $2q+1$, then
\[
k_{1}=1,\quad k_{2}=q, \quad k_{3}=p-2,\quad k_{4}=p-1-q,
\]
are four different integers such that $k_{i}!\equiv 1 \pmod{p}$ for $1\leq i \leq 4$.
\end{lem}

\begin{proof}
Since $p\backslash(q!-1)$, then $p>q$. Otherwise, $q!\equiv 0 \pmod{p}$, which contradicts
that $q!\equiv 1 \pmod{p}$. We already have $1!,q!\equiv 1\pmod{p}$, and by Wilson's symmetry,
$(p-2)!,(p-1-q)!\equiv 1 \pmod{p}$ because $q$ is odd. If the $k_{i}$ were not distinct, then $p-2=1$ if and only if $p=3$, $p-1-q=1$ if and only if $p=q+2$,
$q=p-2$ if and only if $p=q+2$, or $q=p-1-q$ if and only if $p=2q+1$, or $p-2=p-1-q$ if
and only if $q=1$, but all of these are blocked by the hypotheses.
\end{proof}

\section{Main Theorem}

\begin{thm} There exist infinitely many primes $p$ such that
\[
|\{1\leq n\leq p-1:n!\equiv 1 \pmod{p}\}| \geq 4.
\]
\end{thm}
\begin{proof}
Let $q=6t+1$ for an integer $t\geq 1$. It follows that $q$ is odd, while
\[
q+2 = 6t+3 = 3(2t+1)\quad\text{and}\quad 2q+1 = 12t+3 = 3(4t+1),
\]
which implies $q+2$ and $2q+1$ are composite. Thus any prime divisor $p$ of $q!-1$ cannot
be equal to $q+2$ or $2q+1$. Now $q!-1>1$, and hence such a prime $p$ exists. By the valid tetrad
lemma, there exist $k_{1},k_{2},k_{3},k_{4}$, all distinct, such that their factorials
coincide in having residue $1$ modulo $p$.

Let $S$ be a finite set of positive primes with such a property. Choose $q=6t+1$ larger
than any of those primes. Then $q!-1$ has a positive prime divisor $p$ that is larger than $q$, and thus
it cannot belong to $S$, and the result follows.
\end{proof}

\begin{remark}
This result is connected to problem A15 of Guy's compendium \cite[p. 54]{rK04} and the corresponding observation by Mąkowski \cite{Makowski1983}.
\end{remark}

\begin{remark}
For $q=aK+b$, with $a$ even and $b$ odd, the construction of the proof of the main theorem (cf. \cite[Remark 2.8, p. 556]{HS02}) works whenever
\[
\gcd(a,b+2)>1,\quad\text{and}\quad \gcd(a,2b+1)>1,
\]

Indeed, these conditions force $q+2$ and $2q+1$ to be composite for every $K$, so the two possible Wilson-pair collisions are ruled out. The choice $q=6K+1$ is a minimal case and makes the aforementioned pair divisible by same prime divisor, namely $3$.
\end{remark}

\section{Some Computations}

\subsection{Euclidean Scheme}
We can obtain sequences of primes with at least a tetrad of factorials congruent to $1$ using the
Euclidean scheme of the proof (see \cite[Section 4.3]{GKP94}). If we begin
with $t=1$, then we have $q=6\cdot 1+1 = 7$. But $7!-1=5039$ is prime, and the smallest $q$ of the
form $6t+1$ that is larger than $5039$ is $5041$. The least prime factor of $5041!-1$ is $20549$. This was
found considering that if we write the prime divisor as $s=5041+r$, then by Wilson's theorem
\[
 (s-1)!\equiv -1\pmod{s}.
\]

By the definition of $s$, we have
\[
(s-1)! = 5041!5042\cdot 5043 \cdot \cdots \cdot (s-1),
\]
and modulo $s$ we have
\begin{align*}
5042\cdot 5043 \cdot\cdots \cdots (s-1) &\equiv (-r+1)(-r+2)\cdots(-1)\\
&\equiv -(r-1)!\pmod{s}
\end{align*}
because $r$ is even. Thus
\[
 -1 \equiv 5041!(-(r-1)!) \pmod{s}.
\]

Hence $5041!\equiv 1 \pmod{s}$ if, and only if, $(r-1)!\equiv 1 \pmod{s}$. A direct check using the
primes greater than $5041$ with a Python script yields $r=15508$. The next $q$ is $20551$, and the script yields $r=4296$ for the prime factor $s=24847$. After this the algorithm takes more than a few seconds to finish the computation.

If we begin with $t=2$, then $q=6\cdot 2 + 1 = 13$, and the least prime factor of $13!-1$ is $1733$. The
next $q$ of the form $6t+1$ is $1735$. After this the algorithm takes more than a few seconds to finish the computation.

If we begin with $t=3$, then $q=6\cdot 3 +1 = 19$, and the least prime factor of $19!-1$ is $653$. Now the next $q$ of the form $6t+1$ is $655$, and the least prime factor of $655!-1$ is $5407$. After this the algorithm takes more than a few seconds to finish the computation.

\subsection{Tetrads-Only Primes}

It is interesting to calculate the sequence of primes $p$ that admit exactly one tetrad of factorials congruent to $1$ modulo $p$. Using a simple Python script, we can readily find its first terms:
\[
17, 53, 59, 73, 83, 89, 97, 139, 239, \ldots
\]


\section*{Acknowledgements}

The exploration to find the valid tetrad lemma and the main theorem was done with ChatGPT
5.5 Pro. The write up was done by the authors and we take full responsibility for the manuscript.

\section*{Source Code}
\subsection{Euclidean Scheme}
\begin{verbatim}
import sympy as sp

def search_prime_divisor(n):
  # Here n is assumed to be
  # an odd number.
  for s in sp.primerange(n + 1, 10**7):
    r = s - n
    prod = 1
    for u in range(1,r):
      prod = (prod * u) % s
    if prod == 1:
        return s,r
\end{verbatim}

\subsection{Tetrads-Only Primes}

\begin{verbatim}
import sympy as sp

for p in sp.primerange(2, 10000):
    restos = []
    fact = 1

    for n in range(1, p):
        fact = (fact * n) % p

        if fact == 1:
            restos.append(n)

            # Optional early stop:
            # we only care about exactly 4
            if len(restos) > 4:
                break

    if len(restos) == 4:
        print(p)
\end{verbatim}
%\bibliographystyle{amsplain}
%\bibliography{infinite_tetrads}
\printbibliography
\end{document}
